Welche Aufgaben hat der Verantwortliche Aktuar?

ür Lebensversicherer bestimmt VAG § 141, dass jedes Unternehmen einen Ver-
antwortlichen Aktuar zu bestellen hat. Der Verantwortliche Aktuar muss zuverläs-
sig und fachlich geeignet sein. Fachliche Eignung setzt ausreichende Kenntnisse in
der Versicherungsmathematik und Berufserfahrung voraus. Eine ausreichende Be-
rufserfahrung ist regelmäßig anzunehmen, wenn eine mindestens dreijährige Tä-
tigkeit als Versicherungsmathematiker nachgewiesen wird

Zu den wesentlichen Aufgaben des Verantwortlichen Aktuars in der Lebensversi-
cherung gehören (VAG §141 Abs. 5):
-
die Sicherstellung, dass bei der Berechnung der Prämien und der Deckungs-
rückstellungen die rechtlichen Bestimmungen (insb. HGB, DeckRV) eingehalten
werden.
-
die Überprüfung der Finanzlage des Unternehmens daraufhin, ob die dauernde
Erfüllbarkeit der sich aus den Versicherungsverträgen ergebenden Verpflich-
tungen jederzeit gewährleistet ist.

Der Verantwortliche Aktuar vermerkt die versicherungsmathematische Bestäti-
gung unter der HGB-Bilanz, berichtet und er erläutert dem welche Kalkulationsan-
sätze und weiteren Annahmen der Bestätigung zugrunde liegen. Für Versiche-
rungsverträge mit Anspruch auf Überschussbeteiligung legt der Verantwortliche
Aktuar dem Vorstand Vorschläge für eine angemessene Überschussbeteiligung
vor.

Versicherungsunternehmen, die die substitutive Krankenversicherung betreiben,
haben ebenfalls einen Verantwortlichen Aktuar zu bestellen, dessen Aufgaben je-
nen in der Lebensversicherung entsprechen


Erläutern Sie Aufgaben des unabhängigen Treuhänders.

Der Treuhänder nimmt aufsichtsrecht-
liche Befugnisse im Interesse der Versicherten war. Treuhänder zeichnen sich
grundsätzlich durch ihre Zuverlässigkeit, fachliche Eignung und Unabhängigkeit,
wobei letztere insbesondere dann angenommen werden kann, wenn kein Anstel-
lungs- oder sonstiges Dienstverhältnis mit dem Unternehmen oder einem verbun-
denen Unternehmen besteht.

Treuhänder sind im Deutschland für unterschiedliche Bereiche zu bestellen:
-Treuhänder für das Sicherungsvermögen,
-Treuhänder für Prämienänderungen in der Lebens- und Krankenversicherung,
-Treuhänder für Bedingungsanpassungen in der Lebens- und Krankenversiche-
rung.

VAG § 128 fordert die Bestellung eines Treuhänders und eines Stellvertreters zur
Überwachung des Sicherungsvermögens in der Lebens- und in der Krankenver-
sicherung sowie in der privaten Pflegepflicht- und in der Unfallversicherung mit
Prämienrückgewähr.

Zentrale Aufgabe des Treuhänders für das Sicherungsvermögen ist die Bestätigung
im Jahresabschluss unter der Bilanz, dass das Sicherungsvermögen vorschriftsmä-
ßig angelegt und aufbewahrt ist. Bei Streitigkeiten zwischen dem Treuhänder und
dem Versicherer entscheidet die Aufsichtsbehörde.
Prämienänderungen für bestehende Verträge in der Lebens-, Kranken- und Un-
fallversicherung mit Beitragsrückgewähr sowie der Berufsunfähigkeitsversicherung
bedürfen der Zustimmung eines unabhängigen Treuhänders (VAG § 155).

In der Lebensversicherung und in der Krankenversicherung nach Art der Lebens-
versicherung können die Vertragsbedingungen unter bestimmten Bedingungen
mit Zustimmung eines Treuhänders (juristischer Treuhänder) angepasst werden.


Welche Bedeutung haben die sog. „Schlüsselfunktionen“ als Teil des
Governance-Systems der Versicherungsunternehmen?
Unterscheiden Sie die Schlüsselfunktionen hinsichtlich der verfolgten
Ziele und Aufgaben.

Die Geschäftsorganisation schließt zwingend vier Schlüsselfunktionen ein:
-die unabhängige Risikocontrollingfunktion (§ 26 VAG),
-die Compliance-Funktion (§ 29 VAG),
-die Funktion der internen Revision (§ 30 VAG) und
-die versicherungsmathematische Funktion (§ 31 VAG) ein.

Der Begriff „Funktion“ bezeichnet dabei die interne Kapazität innerhalb der Geschäftsor-
ganisation zur Übernahme praktischer Aufgaben.
Darüber hinaus ist sicher zu stellen, dass jede Schlüsselfunktion frei von Einflüssen
ist, die eine objektive und unabhängige Aufgabenwahrnehmung verhindern würde.


Das Risikomanagement verfolgt als oberstes Ziel den verantwortungsvollen Um-
gang mit Risiken, damit die Unternehmensfortführung nicht gefährdet und die dau-
erhafte Erfüllbarkeit der Verpflichtungen gegenüber den Kunden sichergestellt
wird.

Aufgaben der unabhängigen Risikocontrollingfunktion sind insb.:
- die Beurteilung der Konzeption und der Wirksamkeit des Risikomanagementsystems,
-über die Ergebnisse an das Management zu berichten (inkl. ORSA).

Compliance umschreibt die Sicherstellung der Einhaltung von Recht, Gesetz und
interner Regelungen bzw. -vorgaben durch Maßnahmen, wie z.B. Richtlinien, Schu-
lungen und Kontrollen.
die Identifikation,
die Beurteilung,
die Überwachung,
die Vermeidung bzw. Verringerung sowie für
die Berichterstattung bezüglich der Compliance-Risiken der Gesellschaft.

Die Interne Revision ist für die Überprüfung der Funktionalität, der Effektivität
und der Effizienz des internen Kontrollsystems (IKS) sowie allen anderen Be-
standteile des Governance-Systems zuständig. Die Interne Revision muss objektiv
und unabhängig sein, inbs. unterliegt sie nicht den Weisungen des Managements.

Ziel des IKS ist
- die Sicherstellung der Wirksamkeit und die Wirtschaftlichkeit der Geschäfts-
tätigkeit,
- die Einhaltung der Ordnungsmäßigkeit und Verlässlichkeit der internen und
externen Rechnungslegung sowie
- die Einhaltung der für das Unternehmen maßgeblichen rechtlichen Vorschrif-
ten.

Die Versicherungsmathematische Funktion verantworte
- die Koordination der Berechnung der versicherungstechnischen Rückstellun-
gen 
- Entwicklung von Risikomodellen,
-Würdigung aller relevanten Risiken,


Unterscheiden Sie Formalziel und Sachziel der Versicherungsunterneh-
mung und erläutern sie deren spezifische Beziehung.

Während das Formalziel der Unternehmung das handlungsleitende Motiv des Managements fokussiert 
(Bedarfsdeckung, Sicherheit, Wachstum, Gewinn), kennzeichnet das Sachziel die Branche und das 
Geschäftsmodell der jeweiligen Unternehmung (hier: Bereitstellung von Versiche-rungsschutz).

Grundsätzlich besteht ein Zusammenhang zwischen Formal- und Sachziel derart,
dass die Erreichung der formalen Unternehmensziele abhängig von der Qualität
der Sachzielerreichung ist, also beispielsweise der Wettbewerbsfähigkeit der Un-
ternehmung und der Qualität seiner Produkte bzw. Dienstleistungen. In diesem
Sinne hat die Sachzielerreichung grundsätzlich eine gegenüber der Formalzieler-
reichung dienende Funktion.


Was versteht man unter dem Safety-First-Prinzip und welche Bedeutung
hat es für die Steuerung der Versicherungsunternehmung.

Im Zusammenhang mit der Formalzieldiskussion in der Versicherung wird regel-
mäßig der Konflikt zwischen Gewinnerzielungsabsicht und Sicherheitsstreben be-
tont. Unter formalen Aspekten erweisen sich Gewinn- und Sicherheitsstreben hin-
gegen als zwei unterschiedliche Sichtweisen auf eine einheitliche Grundgröße, den
zukünftigen Unternehmenserfolg. Jede Gewinnerzielung, sei es im technischen o-
der nicht-technischen Geschäft, unterliegt wesentlichen, tatsächlichen Risiken.

Insoweit tritt an die Stelle der ökonomischen Optimierung das fortge-
setzte betriebswirtschaftliche Streben der Unternehmung nach Erfolgssteigerung.
M. a. W. gilt es den Unternehmenserfolg unter Beachtung einer kontrollierten Risi-
koexposition zu erhöhen. Dies wird im Zusammenhang mit der Steuerung von Ver-
sicherungsunternehmen als Safety-First-Konzept bezeichnet.

INSERT SLIDE 56 (KPMG)



Problematisieren Sie den „erwarteten Gewinn“ als finanzielle Zielgröße
zur betriebswirtschaftlichen Steuerung der Versicherungsunternehmung.

Theoretisch entspricht der Erwartungswert jedoch lediglich der durchschnittlichen 
Realisation eines Zufallsexperimentes.
Die Interpretation des Erwartungswertes als Maß für ein singuläres Ereignis, wie die Prognose des Unternehmenserfol-
ges in einer bestimmten Periode ist hingegen nicht unproblematisch. Bei unterstellter 
symmetrischer Ergebnisverteilung stellt der Erwartungswert die Größe dar,
die mit 50-prozentiger Wahrscheinlichkeit unterschritten aber auch mit 50-prozen-
tiger Wahrscheinlichkeit überschritten wird.


Welche Vorzüge aber auch welche Probleme ergeben sich aus der Verwendung der 
„erwarteten Eigenkapitalrendite“ als Zielgröße?

Absolute Erfolgskennzahlen, wie z.B. der erwartete Periodengewinn sind für Zwe-
cke der Unternehmenssteuerung häufig ungeeignet, weil sie keine Rückschlüsse auf
die Effizienz des internen Ressourceneinsatzes zulassen (Grundsatz der pretialen
Unternehmenssteuerung).

Als Ersatzzielgrößen für den erwarteten Cash Flow kommen dann relative
Erfolgsmaßstäbe als Rentabilitätskennzahlen zur Anwendung, beispielsweise (er-
wartete) Eigenkapitalrendite oder Return on Risk Adjusted Capital (RORAC).


Was verstehen wir unter dem Wert des Unternehmens und welche Be-
deutung kommt ihm im Rahmen der Unternehmenssteuerung zu?

Aus Sicht der Unternehmenseigner konkretisiert sich der Erfolgsmaßstab im öko-
nomischen Wert der Unternehmung, dem sog. Unternehmenswert als Marktwert
des Eigenkapitals.

Zum Zwecke der Unternehmenswertermittlung von Versicherungsunternehmen wird 
regelmäßig die Verwendung des Equity-Verfahrens als
Spezialfall des Discounted Cash Flow (DCF)-Models vorgeschlagen. Dieses Ver-
fahren ermittelt den Marktwert des Eigenkapitals direkt als Barwert der erwarteten
ausschüttungsfähigen Zahlungsmittel (Free Cash Flow), diskontiert mit einem adä-
quaten Kapitalkostensatz.



Unterscheiden Sie klassische Systeme der Aufbauorganisation und erläutern Sie 
potenzielle Vor- und Nachteile.

Im Einliniensystem sind die Stellen sind dann einfach, im Mehrliniensystem
mehrfach unter- bzw. übergeordnet. Die Aufbauorganisation in zahlreichen Versi-
cherungsunternehmen folgt dem Matrixsystem, das in besonderer Weise geeignet erscheint, 
produkt- und kundenbezogene oder funktions- und spartenbezo-
gene Konflikte transparent und entscheidungsfähig zu gestalten. Beispielsweise
kann so über die Produktgestaltung in der Gebäudeversicherung gemeinsam von
der Stelle „Gebäudeversicherung“ und dem „Vertrieb“ entschieden werden.


In Versicherungsunternehmen unterscheiden wir zwischen primär funktions-
bezogener Aufbauorganisation (Spezialisierung nach Sparten) oder primär
produktbezogenem Aufbau (Spezialisierung nach Funktionen und Vertriebswegen)
oder primär kundengruppenbezogener Aufbauorganisation (Spezialisierung nach
Sparten, und Funktionen)




Erläutern Sie das Konzept der Wertschöpfungskette und unterscheiden
Sie Kernleistungs- und Unterstützungsprozesse in Versicherungsunter-
nehmen.

Das Konzept der Wertschöpfungskette nach Porter strukturiert den betriebli-
chen Leistungserstellungsprozess in Kern- und Unterstützungsleistungen, die mit-
einander in vielfältiger Weise verknüpft bzw. verkettet sind. Ausgangshypothese
ist, jedes “Kettenglied” im Leistungserstellungsprozess trägt in spezifischer Weise
zur Wertentstehung im Unternehmen bei. Dieser Beitrag ist sachlich und wertmä-
ßig isolierbar und unterliegt sodann der Disposition des Managements.


Die Leistungserstellung in der deutschen Versicherungswirtschaft ist traditionell
davon geprägt, dass nahezu jeder Anbieter versucht, sämtliche Kernprozesse
(z. B. Produktentwicklung, Vertrieb, Antragsprüfung, Bestandsverwaltung, Regu-
lierung, Kapitalanlage, Rückversicherung) und Unterstützungsprozesse (z. B.
Zahlungsverkehr, Rechnungswesen, Informationstechnik, Personalverwaltung,
Haustechnik) vollständig innerhalb der eigenen Unternehmung zu organisieren und
durchzuführen. Dies hat dazu geführt, dass sich die über viele Jahre gewachsenen
Strukturen und Abläufe in den einzelnen Häusern in hohem Maße als spezifisch,
differenziert und in der Folge kostenintensiv darstellen.

Die Frage der Neuordnung umfasst dann die Gestal-
tung und Bewertung alternativer Handlungsmöglichkeiten im Spektrum von Eigen-
fertigung (Prozessoptimierung) und Fremdbezug (Outsourcing).

Unter ökonomischer Perspektive ergibt sich hierbei eine Vielzahl relevanter Frage-
stellungen. So sind die Versicherer zunächst bemüht, ihre spezifischen Stärken
und Kernkompetenzen zu identifizieren und Strategien zu entwickeln, um diese
zu erhalten bzw. auszubauen. Andererseits wird geprüft, inwieweit durch Koope-
ration und Auslagerung Leistungs- und Kostenvorteile erreichbar sind (Verbund-
effekte, Skaleneffekte).




Welche Aufgaben umfasst die Vertriebsfunktion im Versicherungsunter-
nehmen?

Die Vertriebsfunktion im Versicherungsunternehmen umfasst zahlreiche praxisre-
gelten auch für die Bereitstellung von Versicherungsschutz das Prinzip
„Absatz vor Produktion“ sowie die Einbringung des „externen Faktors“ als Vo-
raussetzung der Leistungserstellung
levante Aufgaben (z. B. Marktforschung, Absatzplanung, Absatzkanalmanage-
ment, Kommunikations- und Servicepolitik).



Skizzieren Sie wesentliche Prozessschritte der Vertrags- und Schaden-
bearbeitung in Versicherungsunternehmen.

(Betrieb)

Unter dem Begriff des Versicherungsbetriebs werden vor allem Leistungsprozesse
wie Risikoprüfung, Annahmeentscheidung und laufende Verwaltung der Versiche-
rungsbestände zusammengefasst. Im Einzelnen erfolgen hier die innerbetriebliche
Begleitung des Absatzprozesses bis zum Vertragsabschluss, die sich anschließende
Erstbearbeitung (Übermittlung Verbraucherinformationen), die Folgebearbeitungen 
(Bestandsverwaltung, Bestandsführung, Obliegenheiten) während der Laufzeit des Versicherungsvertrages 
sowie die Schlussbearbeitung (Auszahlungen bei Überschuss, Löschung der Daten) bei dessen Ende.

(Schaden)

Aus Sicht der Kunden stellt die Schadenbearbeitung der Schadenabteilung (im
Bereich der Personenversicherung ist die Bezeichnung Leistungsabteilung ge-
bräuchlich) in der Regel den einzigen Anknüpfungspunkt zur Beurteilung der Qua-
lität des Versicherungsschutzes dar.

MORE DETAILS IN P. 77


UNTERSTUEZUNGSPROZESSE EIGENTLICH NUR NENNEN  

Nennen und erläutern Sie die wesentlichen Instrumente des Marketing
Mix in Versicherungsunternehmen

NICHT RELEVANT



Begründen Sie die Charakterisierung der Kapitalanlagetätigkeit als
„ökonomisches Kuppelprodukt“ im Rahmen der Bereitstellung von Ver-
sicherungsschutz.

Unter betriebswirtschaftlichen Gesichtspunkten erzwingt das zeitliche Auseinan-
derfallen von Prämieneinzahlungen und wesentlichen Auszahlungen die Kapitalan-
lagetätigkeit der Versicherer. Dieser auch als wirtschaftliche Kuppelproduktion be-
zeichnete Zusammenhang begründet sich aus der Vorauszahlung der Prämien zum
Vertragsabschlusszeitpunkt, der (stochastischen) Auszahlung im Schadenfall wäh-
rend der Vertragslaufzeit.


Welche grundsätzlichen Einflüsse hat die Kapitalanlagetätigkeit auf die
Erfolgs- und Risikosituation des Versicherungsunternehmens?



Welche Faktoren begründen bzw. begünstigen Konzentrationstendenzen 
in der Versicherungswirtschaft?


Die deutsche Versicherungswirtschaft sieht sich einem verstärkten Wettbewerb
um Kunden, Investoren und Mitarbeiter ausgesetzt. Traditionell erfolgte der
kundenbezogene Wettbewerb unter den deutschen Versicherungsunternehmen
überwiegend auf dem Felde der Distributionspolitik (Vertriebssteuerung) und der
Kommunikationspolitik (Werbung, Öffentlichkeitsarbeit). Mit dem Wegfall der
präventiven Produktkontrolle zum 1. Juli 1994 entstand die Möglichkeit einer weit-
gehend freien Produktpolitik- und Prämiengestaltung.


Deregulierung und Marktöffnung führen zu qualitativen und quantitativen Verän-
derungen im Wettbewerbsumfeld der Versicherer. Konzentrationsstrategien in-
nerhalb der Branche zielen auf betriebswirtschaftliche und versicherungstechni-
sche Wettbewerbsvorteile und damit letztlich auf die Steigerung des Wertes der
Unternehmung selbst und ihrer Leistungsversprechen.



Unterscheiden Sie unterschiedliche Formen der Kooperation zwischen VU's.

INSERT IMAGE 110 (DAV)

Diese Entwicklungen gehen notwendigerweise mit einer verstärkten Wettbe-
werbskontrolle einher. Mit Blick auf die Versicherungswirtschaft besteht hier ein
Spannungsfeld zwischen zulässigen Kooperationen einerseits und nicht-zuläs-
sigem abgestimmten Verhalten oder Ausnutzen einer marktbeherrschenden
Stellung andererseits.

Beispielhaft
genannt seien die Kooperation auf Branchenebene zur Gewinnung einer statistisch
aussagekräftigen Datenbasis für Tarifbildung und Reservierung mit der Folge einer
höheren Kalkulationszuverlässigkeit und risikoadäquaten Prämien. Außerdem exis-
tieren Kooperationen zwischen Versicherungsgruppen in Form von Mitversiche-
rung und Versicherungspools zur Tragung besonders gefährlicher und beson-
ders großer Risiken.
